\documentclass[10pt]{article}

% —— Core packages ——
\usepackage{amsmath,amssymb}
\usepackage{stmaryrd}
\usepackage[utf8]{inputenc}
\usepackage{textcomp}
\usepackage{times}
\usepackage{microtype}
\emergencystretch=1em
\usepackage[most]{tcolorbox}
\tcbuselibrary{theorems,skins,breakable}

% —— Page layout ——
% Compact body, still preserving right margin for handwritten notes.
\usepackage[
  a4paper,
  top=14mm,
  bottom=14mm,
  left=12mm,
  right=62mm
]{geometry}

% Extra header width so the title/rule spans into the note area.
% = right margin (62mm) - left margin (12mm) = 50mm.
\newlength{\notesmargin}
\setlength{\notesmargin}{50mm}

% —— Modal logic and Greek ——
\DeclareUnicodeCharacter{03B9}{\ensuremath{\iota}}
\DeclareUnicodeCharacter{25A1}{\ensuremath{\Box}}
\DeclareUnicodeCharacter{25C7}{\ensuremath{\Diamond}}
\DeclareUnicodeCharacter{03BB}{\ensuremath{\lambda}}
\DeclareUnicodeCharacter{03C6}{\ensuremath{\varphi}}
\DeclareUnicodeCharacter{00AC}{\ensuremath{\neg}}
\DeclareUnicodeCharacter{2227}{\ensuremath{\wedge}}
\DeclareUnicodeCharacter{22A2}{\ensuremath{\vdash}}

\newenvironment{definitionbox}[1][]{%
  \begin{center}
  \begin{tcolorbox}[
    enhanced,
    breakable,
    width=0.9\textwidth,
    colback=gray!2,
    colframe=gray!60,
    coltitle=black,
    fonttitle=\bfseries\small,
    title={Definition: #1},
    boxed title style={
      colback=gray!3,
      colframe=black!60,
      sharp corners,
      boxrule=0.4pt,
    },
    boxrule=0.4pt,
    sharp corners,
    left=4pt,right=4pt,top=4pt,bottom=4pt,
    before skip=6pt,
    after skip=6pt,
  ]}%
  {\end{tcolorbox}\end{center}}

% —— Tables ——
\usepackage{array}
\usepackage{longtable}
\usepackage{booktabs}
\usepackage{calc}
\newcolumntype{C}[1]{>{\raggedright\arraybackslash}p{\dimexpr#1\linewidth-4\tabcolsep\relax}}

% —— Paragraph spacing ——
\usepackage{parskip}
\setlength{\parindent}{0pt}
\setlength{\parskip}{0.25em}

% —— Section formatting ——
\usepackage{titlesec}
\usepackage{needspace}
\usepackage{etoolbox}

\setcounter{secnumdepth}{3}

\titleformat{\section}[hang]
  {\normalfont\large\bfseries}
  {\thesection.}{0.7em}{}

\titleformat{\subsection}[hang]
  {\normalfont\normalsize\bfseries}
  {\thesubsection.}{0.7em}{}

\titleformat{\subsubsection}[hang]
  {\normalfont\small\bfseries}
  {\thesubsubsection.}{0.7em}{}

\titlespacing*{\section}{0pt}{1.0em}{0.35em}
\titlespacing*{\subsection}{0pt}{0.75em}{0.25em}
\titlespacing*{\subsubsection}{0pt}{0.55em}{0.15em}

% —— Optional page break before each section except first ——
% Disabled for compactness. Uncomment if you want each section on a new page.
% \newif\iffirstsection
% \firstsectiontrue
% \pretocmd{\section}{%
%   \iffirstsection\firstsectionfalse\else\clearpage\fi
% }{}{}

% —— Lists ——
\usepackage{enumitem}
\setlist{
  nosep,
  leftmargin=1.5em,
  topsep=0.15em,
  partopsep=0pt,
  parsep=0pt,
  itemsep=0.05em
}
\setlist[itemize,1]{label=\textbullet}
\setlist[itemize,2]{label=\textendash}

% —— Figures ——
\usepackage{graphicx}
\usepackage{float}
\setkeys{Gin}{keepaspectratio,width=0.65\linewidth}
\makeatletter
\def\fps@figure{H}
\makeatother
\graphicspath{{figures/}}

% —— URL/DOI wrapping ——
\usepackage{xurl}

% —— CSL bibliography ——
\newlength{\cslhangindent}
\setlength{\cslhangindent}{1.2em}
\newenvironment{CSLReferences}[2]%
  {\begin{enumerate}[label={[\arabic*]},leftmargin=\cslhangindent,itemsep=0pt,parsep=0pt,topsep=0pt]%
   \raggedright\sloppy
   \renewcommand{\bibitem}[2][]{\item}}%
  {\end{enumerate}}
\providecommand{\CSLLeftMargin}[1]{\item[#1]}
\providecommand{\CSLRightInline}[1]{#1}

% —— Hyperref ——
\usepackage[colorlinks=true, linkcolor=black, urlcolor=blue, citecolor=blue]{hyperref}
\urlstyle{same}

% —— Fix for Pandoc/longtable "none" counter ——
\makeatletter
\@ifundefined{c@none}{%
  \newcounter{none}%
}{}
\makeatother
\renewcommand{\thenone}{}
\providecommand*{\theHnone}{\arabic{none}}

% —— Line spacing and general flow ——
\linespread{1.08}
\raggedbottom

% —— Block quotes ——
\renewenvironment{quote}{%
  \par\vskip 0.25em
  \leftskip=1.5em\rightskip=1.5em\small\itshape
  \noindent\ignorespaces
}{\par\vskip 0.25em}

% —— Pandoc compatibility ——
\providecommand{\tightlist}{}
\providecommand{\pandocbounded}[1]{#1}

% —— Citeproc fixes ——
\makeatletter
\providecommand{\citeproctext}{}
\let\@oldprotect\protect
\def\@citeprocprotection{\let\protect\@oldprotect}
\newcommand{\citeproc}[2]{#2}
\protected\def\citeprocfootnote#1{\footnote{#1}}
\makeatother

% —— Special character definitions ——
\DeclareUnicodeCharacter{2203}{\ensuremath{\exists}}
\DeclareUnicodeCharacter{2200}{\ensuremath{\forall}}
\DeclareUnicodeCharacter{2228}{\ensuremath{\vee}}
\DeclareUnicodeCharacter{2192}{\ensuremath{\rightarrow}}
\DeclareUnicodeCharacter{2194}{\ensuremath{\leftrightarrow}}
\DeclareUnicodeCharacter{2234}{\ensuremath{\therefore}}
\DeclareUnicodeCharacter{2248}{\ensuremath{\approx}}
\DeclareUnicodeCharacter{2208}{\ensuremath{\in}}
\DeclareUnicodeCharacter{2260}{\ensuremath{\neq}}
\DeclareUnicodeCharacter{2713}{\checkmark}
\DeclareUnicodeCharacter{2212}{-}

% —— Angle brackets for LPCs ——
\newcommand{\llangle}{\langle\!\langle}
\newcommand{\rrangle}{\rangle\!\rangle}
\DeclareUnicodeCharacter{27EA}{\ensuremath{\llangle}}
\DeclareUnicodeCharacter{27EB}{\ensuremath{\rrangle}}

% —— Custom commands for frequent notation ——
\newcommand{\Rfr}[2]{\ensuremath{\mathcal{R}(#1, #2)}}
\newcommand{\LPC}[1]{\ensuremath{\llangle#1\rrangle}}

% —— Rule colour ——
\usepackage[dvipsnames]{xcolor}
\definecolor{ruleColor}{gray}{0.60}

% ============================================================
\begin{document}
% ============================================================

% —— Header block: spans body + notes margin ——
\noindent\begin{minipage}{\dimexpr\textwidth+\notesmargin\relax}

  \noindent{\sffamily\Large\bfseries Chapter 3 --- From Russell to Kripke: The Argument for Rigid Designation}%
  \hfill{\sffamily\normalsize 2026-04-22}

    \vspace{2pt}
  \noindent{\sffamily\small\itshape Revised draft, aligned with Ch. 2 originality audit}
  
  \vspace{4pt}
  \noindent{\color{ruleColor}\rule{\linewidth}{0.35pt}}

\end{minipage}

\vspace{3pt}

% —— Body: compact left column, wide right margin for notes ——
\section{Chapter 3 --- From Russell to Kripke: The Argument for Rigid Designation}\label{chapter-3-from-russell-to-kripke-the-argument-for-rigid-designation}

Chapter 2 closed at the point where the necessity-of-identity formula, applied to names, presupposes that names exhibit the same reference-invariance across worlds that variables have under assignment. That presupposition is a substantive metasemantic claim; the formula does not establish it. This chapter presents the independent argument Kripke makes for it in \emph{Naming and Necessity}, lays out the received view on why descriptions cannot supply what the argument requires, and reconstructs a dialectical reading on which Russell first identified the mechanism and then declined to generalize it. The chapter closes with a synthesis of the leading cognitive implementations --- Evans, Recanati, Devitt --- on which each identifies a genuine desideratum but shares a representational template that stops short of the primitive tracking capacity the thesis requires. The gap this reading identifies is the one Chapters 4 and 5 undertake to close.

The chapter does not claim to refute rigid designation. The semantic thesis, as Kripke states it, is taken as well-motivated and largely correct. The three sections that draw on the literature --- §3.1 (Kripke's arguments), §3.2 (the scope analysis), and parts of §3.4 (baptism and the causal chain) --- present the post-Kripkean consensus and are not claimed as original. The modest positive contribution is concentrated in §3.3 (the reading of Russell's passage), §3.5 (the unified reading of the three gaps), and §3.6 (the synthesis of the cognitive implementations). Each of these is flagged as a reading rather than a result. Kripke is explicit that a full account of reference would supply a cognitive mechanism and that he is not supplying it; the chapter takes him at his word and points to where the rest of the dissertation picks up.

\section{3.1 Rigid Designation: Kripke's Independent Argument}\label{rigid-designation-kripkes-independent-argument}

This section presents the post-Kripkean consensus on the independent argument for rigid designation. No analytical move in §3.1 is claimed as original; the section's purpose is to put the case for the thesis cleanly enough that later sections can refer back to it.

Kripke's thesis is that proper names are rigid designators: a proper name designates the same individual in every possible world in which that individual exists.\footnote{Kripke, \emph{Naming and Necessity} (1980), p.~48.} The thesis is defended in \emph{Naming and Necessity} by three converging arguments, none of which depends on the necessity-of-identity formula of Chapter 2. The formula regiments the modal consequence once rigidity is granted; the case for rigidity runs separately, and is now canonical.\footnote{The three arguments are conventionally grouped under these labels in the teaching literature. See Speaks, ``Kripke on the necessary a posteriori'' (lecture notes, Notre Dame PHIL 83104, 2011); LaPorte, ``Rigid Designators,'' \emph{Stanford Encyclopedia of Philosophy}, §1; Salmon, \emph{Reference and Essence} (1981), and Stanley, ``Names and Rigid Designation'' (1997).}

The modal argument turns on counterfactual intuition. Aristotle could have died in infancy, never taught Alexander, never written the \emph{Nicomachean Ethics}. In that counterfactual scenario, it is Aristotle who failed to do these things, not someone else. The name ``Aristotle'' tracks the individual across counterfactual variation, even when every description we associate with him fails. If ``Aristotle'' meant ``the teacher of Alexander who wrote the \emph{Nicomachean Ethics},'' then in a world where no one satisfied that description the name would refer to no one, or, under a descriptivist semantics, to whoever happened to satisfy the description. Both readings clash with the intuition that Aristotle is the same man across counterfactuals. The intuition is the test: the modal sentence ``Aristotle might have died in infancy'' is true, and its truth fixes the referent across the counterfactual.

The epistemic argument turns on ignorance and error. A competent user of ``Feynman'' may know only that Feynman was a physicist, or may believe things about him that are factually wrong. On a descriptivist theory such a user either fails to refer or refers to whoever fits her associated description, which in cases of error will be someone other than Feynman. Kripke's observation is that such speakers manage to refer to Feynman regardless. The sharpest version of the case is Gödel/Schmidt: suppose, counterfactually, that the proof we attribute to Gödel was actually due to Schmidt, and Gödel merely found and published it. On a descriptivist theory built around ``the prover of the incompleteness theorem,'' ``Gödel'' would refer to Schmidt --- the actual prover. Kripke's intuition is that ``Gödel'' still refers to Gödel.\footnote{Kripke, \emph{Naming and Necessity}, pp.~83--87.} Reference runs through the name itself, not through the descriptions attached to it.

The semantic argument targets substitutivity in modal contexts. ``Necessarily, Benjamin Franklin is Benjamin Franklin'' is true; ``Necessarily, the inventor of bifocals is the inventor of bifocals'' is also true; but ``Necessarily, Benjamin Franklin is the inventor of bifocals'' is false --- he might never have invented them. The name and the description share an actual-world referent but differ in modal profile. They cannot be synonyms. The same point is sharpened by the Hesperus/Phosphorus material later in \emph{Naming and Necessity}: a true identity statement between two rigid designators is necessary, even when the identity was discovered empirically.\footnote{Kripke, \emph{Naming and Necessity}, pp.~100--105.}

These three arguments, together, motivate the distinction between \emph{de jure} rigidity (characteristic of proper names, rigid by stipulation) and \emph{de facto} rigidity (attaching to descriptions whose satisfiers are necessarily the same across worlds, such as ``the least prime'').\footnote{Kripke, \emph{Naming and Necessity}, p.~21, fn. 21.} Names are rigid by how they function as designators, not by what they happen to pick out. The standard modal test that emerges from this material --- used throughout the rest of the chapter --- is the substitution test: a designator \(d\) is rigid if, in every modal context where \(d\) refers, \(d\) refers to the same individual across the worlds the context quantifies over. The thesis and its defence are Kripke's; the presentation in this section follows the standard reconstruction in Speaks's lecture notes, the Stanford Encyclopedia, and the cluster-theory critique literature.\footnote{Kripke, ``Identity and Necessity'' (1971), p.~144, quoting Russell.}

\section{3.2 Why Descriptions Cannot Supply Rigidity}\label{why-descriptions-cannot-supply-rigidity}

This section presents the received view on why descriptions, by themselves, cannot supply the kind of designation Kripke's argument requires. The scope analysis and the \emph{dthat}-style response, together with Kripke's reply to them, are drawn from Recanati 1993, Soames 2002, and Kripke's own preface; none of the analytical moves in §3.2 is claimed as original.\footnote{Recanati, \emph{Direct Reference} (1993), chs.~1--3; Soames, \emph{Beyond Rigidity} (2002), esp.~chs.~2--3; Kripke, \emph{Naming and Necessity}, preface and p.~12 fn. 15.}

Descriptions pick out satisfiers world by world. ``The inventor of bifocals'' picks out whoever, in a given world, invented bifocals. Different worlds can have different inventors; the description tracks whichever individual happens to satisfy it. This world-relative behavior is constitutive of descriptions: they specify a condition and the reference is wherever the condition is met. No accumulation of conditions changes this. ``The stout inventor of bifocals who was a Postmaster General'' still picks out a satisfier world by world; adding conditions narrows the range of worlds where the description is satisfied but does not make the reference invariant across them.

The formal expression of this is the scope problem. In a sentence like \(\Box \, F(\iota x\, G x)\), the description \(\iota x\, G x\) can take scope inside or outside the modal operator. Read with narrow scope, ``necessarily, \(F\) of whoever is \(G\)'' requires that in every world the satisfier of \(G\) in that world is also \(F\). Read with wide scope, ``of whoever is actually \(G\), necessarily \(F\)'' requires that the actual-world satisfier of \(G\) is \(F\) in every world. The wide-scope reading behaves modally as if the description were rigid: the referent is fixed in the actual world and held constant. But the rigidity is positional, not constitutive --- it is achieved by choice of scope, not by anything intrinsic to the description.\footnote{Recanati, \emph{Direct Reference} (1993), ch.~1; see also Recanati, ``Rigidity and Direct Reference'' (1988).} Soames develops this contrast at length: a description with wide scope and a rigid designator share modal profiles in simple cases but come apart on more complex constructions; Soames's ``obstinately'' rigid designators (which keep their referent fixed across all worlds, including worlds where the referent does not exist) cannot be modeled by any scope manipulation of a description.\footnote{Soames, \emph{Beyond Rigidity} (2002), pp.~39--58.}

Kaplan's \emph{dthat} operator is standardly read in the same light.\footnote{Kaplan, ``Demonstratives'' (1989), pp.~521--24.} \(dthat(\iota x\, F x)\) takes a description and converts it into a directly referring term whose reference is fixed in the actual world and held constant across worlds. On the consensus reading, \emph{dthat} does the same logical work as wide-scope \(\iota x\, F x\), but it does the work through rigidification rather than scope manipulation --- and, crucially, it makes the rigidification intrinsic to the designator rather than dependent on how the designator is positioned relative to a modal operator. Names differ from rigidified descriptions in a further respect: for names, wide scope is the only reading, because the name has no descriptive structure that could generate a narrow-scope reading in the first place. ``Aristotle'' does not pick out whoever happens to be Aristotle in each world, because ``being Aristotle'' is not a condition the name imposes.

The upshot, on the received view, is that descriptions achieve rigidity only derivatively --- by scope choice in natural-language cases, or by stipulated rigidification via \emph{dthat} in the formal case --- whereas names are rigid independently of scope. Kripke himself makes this distinction in the preface and in the scope footnote: rigidity is not the same as wide-scope behavior, and the rigid-designation thesis is not a thesis about scope.\footnote{Kripke, \emph{Naming and Necessity}, preface (1980 ed., pp.~6--15) and p.~12 fn. 15. Kripke addresses the wide-scope strategy explicitly in his reply to Dummett; the position is that wide-scope behavior is compatible with rigidity but does not constitute it.} Chapter 4 will return to this contrast under a different vocabulary, but for present purposes the descriptive side of the asymmetry is settled literature.

\section{3.3 Russell's Mechanism and His Retreat}\label{russells-mechanism-and-his-retreat}

The previous section showed that descriptions, taken on their own, do not supply rigidity. This raises the question of where rigidity does come from. Kripke's positive answer, baptism and the causal chain, is the subject of §3.4. Before turning to it, this section reads a passage that Kripke himself draws on, in which Russell --- working from a different direction --- identifies what, on the reading developed here, is the cognitive mechanism that Kripke's later picture relies on without spelling out. The reading is mine; it is not a claim about what Russell really meant. It does, however, sit in a tradition: Sainsbury, Pears, and Evans have all read Russell's acquaintance principle as extendable beyond sense data once the epistemological constraint is relaxed,\footnote{Sainsbury, \emph{Russell} (1979), on the scope of acquaintance; see also Pears, ``Russell's Theory of Knowledge'' (1972).} and the move I make below is in that broad spirit.

Kripke quotes Russell in ``Identity and Necessity'':

\begin{quote}
So, Russell concludes, if we want to reserve the term ``name'' for things which really just name an object without describing it, the only real proper names we can have are names of our own immediate sense data, objects of our own immediate acquaintance. The only such names which occur in language are demonstratives like ``this'' and ``that.'' And it is easy to see that this requirement of necessity of identity, understood as exempting identities between names from all imaginable doubt, can indeed be guaranteed only for demonstrative names of immediate sense data; for only in such cases can an identity statement between two different names have a general immunity from Cartesian doubt.\footnote{Kripke, ``Identity and Necessity'' (1971), p.~144, quoting Russell.}
\end{quote}

Two things are in this passage, and on the reading developed here they should be separated.

The first is Russell's identification of a mechanism. Russell isolates a specific kind of designator --- the demonstrative --- and characterizes it by what it does not do: it does not describe. The demonstrative ``just names an object,'' without descriptive intermediary. The reference is established by the demonstrative act itself, not by a description the speaker has in mind. On the reading I propose, this is the mechanism the chapter has been after. The descriptive theories cannot supply rigidity because they cannot supply this; the formal scope manipulations of §3.2 simulate this at the level of the designator without supplying it at the level of what the speaker does. Russell, on this reading, is naming the cognitive operation that the rigid-designation thesis presupposes but does not itself describe.\footnote{The term \emph{local rigidity} is introduced in §4.5 and developed in Chapter 5. Anticipations in the literature are discussed there.}

The second is Russell's retreat. Having identified the mechanism, he restricts it to ``immediate sense data'' --- the subset of demonstrative cases in which the reference is exempt from Cartesian doubt. The restriction is epistemological. Russell requires that the referent of a genuine name be something one cannot be wrong about, and only immediate sense data meet this requirement. An ordinary physical object, as a possible hallucination or misperception, fails the Cartesian test. Under Russell's standard, even ``this chair'' is not a genuine name, because one could in principle be wrong about whether there is a chair there at all.

The dialectical position I will exploit in Chapter 5 can be stated from here. Russell's mechanism --- demonstrative grounding --- is the kind of mechanism the rigid-designation literature needs. But by restricting it to incorrigible cases, Russell makes the mechanism unavailable for ordinary names. If the only rigid designators are names of sense data, and ordinary proper names are not rigid (because not grounded in infallible acquaintance), then Kripke's thesis cannot be Russellian in Russell's own terms. The reading this dissertation will develop is that what Russell offers as the mechanism is the right mechanism, and what he requires of the mechanism --- Cartesian infallibility --- is too strong for it to cover the cases Kripke wants to cover.

This reading is not a claim about what Russell really meant. It is a use of Russell for a purpose Russell would not have endorsed. Sainsbury and others have read Russell's acquaintance principle as extendable beyond sense data once the epistemological constraint is relaxed; Evans's generalization of Russell's Principle in \emph{The Varieties of Reference} is in the same spirit.\footnote{Evans, \emph{The Varieties of Reference} (1982), ch.~4.} The present reading sits in this tradition and does not pretend to be alone in it. Its particular shape --- the distinction between \emph{fallibility} and \emph{instability} that will be developed in Chapter 5 --- may be more specific to the PRU account. Demonstrative grounding, on this reading, does not need to be infallible to do the work the rigid-designation thesis requires; it needs to be stable. Russell collapses fallibility and instability because under the Cartesian standard any fallibility just is an instability. That collapse is what the subsequent chapters drop.

The dialectical position Chapter 5 will occupy is visible from here. Russell found the mechanism. Kripke extended the phenomenon (rigid designation covers ordinary names, not just sense-data demonstratives) but left the mechanism unspecified. The gap between them is what needs closing. Closing it requires the mechanism Russell identified, applied at the scope Kripke extended --- which means dropping Russell's Cartesian requirement. Demonstrative grounding does not need to be infallible to do the work. It needs to be stable. The distinction between fallibility and instability is the distinction Russell collapses. Chapter 5 returns to this.

\section{3.4 Baptism and the Causal Chain}\label{baptism-and-the-causal-chain}

This section presents Kripke's positive picture of how names attach to their referents. The exposition follows Kripke's text closely; the reading-fixing vs.~meaning distinction added below is also Kripke's, not original to the chapter.

Kripke's positive account of how names attach to their referents runs through baptism and causal-historical transmission. At some initial act --- the baptism --- a name is attached to its bearer, either by ostension (``let us call this child `Feynman'\,'') or by reference-fixing description (``let `Jack the Ripper' denote whoever committed these murders''). From the baptism, the name propagates through a chain of uses: each speaker acquires the name from an earlier speaker, who acquired it from an earlier speaker, and so on back to the baptism. The speaker at the end of the chain refers to the baptismal referent even if she knows nothing about him beyond the name.\footnote{Kripke, \emph{Naming and Necessity}, pp.~91--97; on reference-fixing vs.~meaning-giving, see pp.~53--60.}

The picture is presented, by Kripke's own explicit framing, as a ``better picture'' rather than as a theory. Kripke is emphatic that he is not offering necessary and sufficient conditions for reference: ``I'm not, in fact, trying to give any sort of necessary and sufficient conditions for reference\ldots{} in any rigorous sense.''\footnote{Kripke, \emph{Naming and Necessity}, p.~94.} The self-characterization matters for the dissertation. Kripke's picture tells us what reference looks like from the outside --- a baptismal act, a chain of transmission, a referent preserved --- but not what cognitive events realize the picture in speakers.

The picture has two parts worth distinguishing. The baptism fixes the reference. The causal chain preserves it across uses, across speakers, across time. Both parts are necessary. A baptism without a chain would fix reference for one speaker, no further; a chain without a baptism would transmit nothing.

Kripke further distinguishes reference-fixing from meaning-giving.\footnote{Kripke, \emph{Naming and Necessity}, pp.~53--60.} A description used in baptism --- ``let `Jack the Ripper' denote whoever committed these murders'' --- fixes the reference of the name but does not thereby become the name's meaning. The name, once fixed, is rigid; the description used to fix it is not. A speaker who later learns that the Ripper was not the person associated with the canonical descriptions does not thereby change the reference of the name. This distinction matters for the present chapter because it clarifies what the causal chain transmits: the name, with its fixed reference, not the reference-fixing description. What is inherited down the chain is the designator and its rigidity, not the descriptive route by which the reference was originally established.

What the picture does well, it does by contrast with descriptivism. Reference does not depend on a speaker's descriptive grasp of the referent. A child can learn ``Aristotle'' from a parent and refer to Aristotle with no descriptive content beyond the name. The picture handles ignorance, error, and the accretion of mistaken belief that actual linguistic practice exhibits.

What the picture leaves open is the cognitive level. That is the subject of §3.5.

\section{3.5 The Explanatory Gap}\label{the-explanatory-gap}

It is worth being explicit about what Kripke's semantic account does not claim to supply. The post-Kripkean literature has widely acknowledged what LaPorte characterizes as the ``promissory note'' structure of the causal-historical picture: it gestures at cognitive and psychological mechanisms it does not itself supply.\footnote{LaPorte, ``Rigid Designators,'' \emph{Stanford Encyclopedia of Philosophy}, §4, characterizing Kripke's appeal to stipulation as ``more like a promissory note than the satisfaction of an explanatory obligation.''} The three gaps noted below are not original diagnoses; they are the standard way of itemizing what the semantic picture leaves open.

First, no account is given of what constitutes the rigid tracking at the cognitive level. Names are said to be rigid; the rigidity is explained semantically as designation of the same individual across worlds; but the cognitive capacity that allows a speaker to use a name as rigid --- to deploy it in modal contexts where descriptions would fail, to apply it in counterfactual scenarios where its referent has none of its actual-world properties --- is not characterized. Kaplan's \emph{dthat}, as we saw, stipulates rigidity at the semantic level; it does not explain how ordinary speakers produce rigidity without a \emph{dthat}-operator in their linguistic toolkit.

Second, no account is given of the cognitive mechanism of baptism. Baptism is described as reference-fixing, but what a speaker does when she baptizes --- what cognitive event supports the ostensive or descriptive attachment of a name to a particular --- is left open. The ostensive case is especially pressing: pointing at a thing and saying ``let us call this `X'\,'' presupposes a capacity to single out \emph{this} thing from the perceptual field, a capacity that descriptive theories also presuppose but do not account for.

Third, no account is given of how the causal chain is psychologically realized. A causal chain of uses requires, at each link, that a speaker learn the name from an earlier speaker and take it up in a way that preserves its reference. What the speaker acquires, cognitively, when she acquires the name --- whether it is a descriptive file, an ability to recognize the bearer, a disposition to accept testimony from certain sources --- is not specified. The chain is characterized at the social-linguistic level; its cognitive substrate is absent.

These three gaps are connected, and the connection is the reading on which the rest of the dissertation turns. On the reading I propose --- and it is a reading, not a result --- the three gaps are three aspects of a single missing element: the primitive tracking capacity that lets a speaker single out a particular and continue to single out that same particular over time, through change, and under modal variation. Rigid tracking, on this reading, is the same capacity exercised at the modal level; baptism is the same capacity exercised in attaching a name to a particular; the causal chain is the same capacity exercised in taking up a name from another speaker and continuing to apply it to that particular. The reading is not mandatory. The three gaps could be filled by three independent accounts: a theory of modal cognition, a theory of ostension, and a theory of testimonial uptake. But the unified reading is explanatorily preferable on the account this dissertation develops, because the primitive it identifies --- local rigidity, the demonstrative tracking Russell named and restricted --- is the same primitive across all three contexts.\footnote{That the cognitive question is real and unresolved is widely recognized; see Dickie, \emph{Fixing Reference} (2015), Introduction and ch.~1.}

\section{3.6 Existing Cognitive Implementations}\label{existing-cognitive-implementations}

The previous section identified the cognitive gap left by the semantic account. Several philosophers have attempted to fill it. This section discusses three: Evans's discriminating-conception account, Recanati's mental-files account, and Devitt's causal-descriptive account. The discussion offers a particular synthesis on which all three share a structural feature that, on the reading developed in this chapter, stops short of the primitive the rest of the dissertation will develop. The individual critiques offered below do not originate here --- Recanati critiques Evans's propositional framing in \emph{Mental Files} (2012), ch.~2; Dickie raises structural worries about file architectures; direct-reference theorists have long charged Devitt with residual descriptivism.\footnote{Recanati critiques Evans's propositional framing in \emph{Mental Files} (2012), ch.~2; Dickie, \emph{Fixing Reference} (2015), raises structural worries about file architectures; on direct-reference critiques of Devitt, see Salmon, \emph{Reference and Essence} (1981), and the discussion in Braun, ``Understanding Belief Reports'' (1998).} The synthesis --- the unified diagnosis that all three share a representational template --- is the section's modest contribution, and is offered as a reading.

Evans's \emph{Varieties of Reference} defended a Russellian principle: genuine singular thought about an object requires that the thinker possess a ``discriminating conception'' of that object --- a way of singling it out that distinguishes it from all other objects.\footnote{Evans, \emph{The Varieties of Reference} (1982), ch.~4.} The discriminating conception is not a description in the cluster-theoretic sense; Evans is careful to distinguish his account from descriptivism. But it is a propositional capacity: to have a discriminating conception of \(a\) is to be able to think, of \(a\), that it is the unique satisfier of some condition the thinker grasps. The strength of Evans's account is its attention to the cognitive level Kripke leaves open. Its limit, on the reading developed here, is the propositional framing: the discriminating conception is given as a kind of thought a thinker can have, rather than as a tracking capacity that supports the kinds of thoughts the thinker has.\footnote{Goodman (2024) surveys the post-Evans trajectory; the question whether Evans's ``conception'' should be read as propositional or as a tracking capacity is pressed in Recanati's response to Evans in \emph{Mental Files}.}

Recanati's \emph{Mental Files} offers a different model.\footnote{Recanati, \emph{Mental Files} (2012).} A speaker's grasp of a singular term is implemented as a ``file'' --- a cognitive structure that collects information about the referent, supports updating as new information comes in, and underwrites the speaker's capacity to think about the referent in absence. Recanati's account improves on Evans's in several respects: files are not propositional in Evans's sense, they support coreference without requiring shared descriptive content, and they have a developmental story that Evans's discriminating conception lacks. The container metaphor, however, retains a representational template. The file is something the speaker has; the referent is what the file is about. The act of opening a file presupposes a capacity to single out a particular and start a file on it. That capacity is not, on Recanati's account, itself a file; it is the precondition for files.

Devitt's \emph{Designation} offers a third approach: causal-descriptivism, which combines Kripke's causal chain with a residual descriptive component to handle Twin-Earth-style cases and cases of reference shift.\footnote{Devitt, \emph{Designation} (1981), ch.~2; see also Devitt and Sterelny, \emph{Language and Reality} (1999), ch.~3.} Devitt's account is closer to Kripke than the previous two in its acceptance of the causal chain as a substantive component; its limit, from the perspective developed here, is that the descriptive component continues to do the heavy lifting in cases where the chain itself is contested. The architecture is still representationalist: reference is a relation between a representational state (the speaker's causal-descriptive complex) and a referent.

The reading I propose is that these three proposals share a structural feature. All three treat reference as a relation between a representational state (a conception, a file, a description-plus-causal-link) and a referent. Each account improves on its predecessors by loosening the descriptive component --- Evans by making the conception a capacity rather than a description, Recanati by making the file a container rather than a set of descriptions, Devitt by supplementing descriptions with causal chains. But the representational template remains. Reference is something the speaker \emph{has} (a conception, a file, a causal-descriptive state); the referent is what the representational state is \emph{about}. What none of the three supplies is the mechanism by which the having is accomplished: the cognitive process through which a speaker's state gets coupled to a particular individual in the first place.

The claim is not that the representational accounts are wrong about what they describe. Evans's discriminating conception is real; Recanati's files exist; Devitt's causal-descriptive hybrids do explanatory work. The claim is that each of them describes the output of a more basic cognitive process --- the tracking capacity Russell identified in demonstrative acquaintance and then set aside as too narrow to generalize. Chapters 4 and 5 develop that tracking capacity under the PRU framework, under the name \emph{local rigidity}, and argue that baptism and the causal chain are best understood as the social stabilization of local rigidity across absence. On the reading offered in this section, the existing cognitive implementations leave a specific absence: the primitive tracking capacity from which reference-relations, as described by Evans, Recanati, and Devitt, are derived.

\section{3.7 Summary and Transition}\label{summary-and-transition}

This chapter has done three things. It has presented Kripke's independent argument for rigid designation --- the modal, epistemic, and semantic arguments, together with the Gödel/Schmidt and related cases --- so that the presupposition Chapter 2 identified now stands as a well-motivated thesis rather than an unargued assumption. It has laid out the received view on why descriptions cannot supply rigidity: the scope problem, the asymmetry with names, and the role of Kaplan's \emph{dthat} as explicit rigidification. And it has offered a particular reading of the Russellian background on which Russell identified the cognitive mechanism (demonstrative grounding) but restricted it by a Cartesian requirement the dissertation will drop.

The chapter's modest positive contribution lies in §3.5 and §3.6: a unified reading of the explanatory gap left by the semantic picture, and a synthesis on which the leading cognitive implementations --- Evans, Recanati, Devitt --- each address the gap but share a representational template that stops short of the primitive tracking capacity the thesis requires. These readings are offered as setup for the positive account, not as discoveries about the literature.

Chapter 4 introduces the analytical vocabulary that makes the primitive visible: the \(x^o\)/\(x^c\) mode distinction, the grounding relation \(\mathcal{R}\), and the notion of local rigidity as the cognitive ground of the Kripkean phenomenon. Chapter 5 uses that vocabulary to develop the positive account: baptism as the lifting of local rigidity into a socially portable name-concept, the causal chain as the testimonial network that sustains the lifting, and PRU as the metasemantic framework that explains how agents come to inhabit Kripkean models without revising the models themselves.

\vfill
\noindent\rule{\textwidth}{0.35pt}\\[1pt]
{\tiny florin.cojocariu@s.unibuc.ro, \today}

\end{document}